% SEGMENT OLP-0481-S01
% Part: applied-modal-logics
% Chapter: temporal-logic
% Section: possible-histories

\documentclass[../../../include/open-logic-section]{subfiles}

\begin{document}

\olfileid{aml}{tl}{poss}

\olsection{சாத்தியமான வரலாறுகள்}

அணுகுறவின் மீது எந்தக் கட்டுப்பாட்டையும் விதிக்க வேண்டியதில்லை என்பதால், நாம்
பயன்படுத்திவந்த காலத் தருக்கத்தின் தொடர்புமுறை மாதிரிகள் மிகவும்
நெகிழ்வானவை. கால மாதிரிகள் கடந்தகாலத்திலும் எதிர்காலத்திலும் கிளைப்படையலாம்
என்பது இதன் பொருள். ஆனால் நிகழ்வுத் தொடர்களைச் சாத்தியமான வரலாறுகளாகக்
கருதும், கிளைப்படைதல் குறித்த இன்னும் ``வாய்ப்புநிலை'' சார்ந்த கருத்தை நாம்
ஆராய விரும்பலாம். இதற்கு நமது மொழியை மாற்றுவது கட்டாயமில்லை. எனினும், நமது
``வழக்கமான'' வாய்ப்புநிலைச் செயற்குறிகள் $\Box$,~$\Diamond$ ஆகியவற்றைச்
சேர்க்கலாம்; காலப்போக்கில் முகவர்களின் அறிவில் ஏற்படும் மாற்றங்களைக் குறிக்க
அறிவுநிலை அணுகுறவுகளைச் சேர்ப்பதையும் ஆராயலாம்.

\begin{defn}
  கால மொழிக்கான ஒரு \emph{சாத்தியமான வரலாறுகளின் மாதிரி} என்பது
  $\mModel{M} = \tuple{T, C, V}$ என்ற மும்மை; இதில்
  \begin{enumerate}
    \item $T$ என்பது கால நிலைகளாகப் பொருள்கொள்ளப்படும் காலியல்லாத கணம்.
    \item $C$ என்பது ஒரு அமைப்பின் கணிப்புப் பாதைகளின், அல்லது
      \emph{சாத்தியமான வரலாறுகளின்}, கணம். வேறு சொற்களில், $C$~என்பது
      $\sigma$ எனும் நிலைத் தொடர்முறைகளின் கணம்; அவை $s_1$, $s_2$,~$s_3$,
      \dots என்ற வடிவில் இருக்கும், மேலும் ஒவ்வொரு $s_i \in T$.
    \item ஒவ்வொரு !!{propositional variable}~$p$ க்கும் காலப் புள்ளிகளின்
      கணம்~$V(p)$ ஐ மதிப்பளிக்கும் சார்பு $V$.
  \end{enumerate}
  எளிமைக்காக, ஒரு வரலாறு~$C$ இல் இருந்தால் அதன் எல்லாப் பின்னொட்டுகளும்
  அக்கணத்தில் இருக்கும் என்றும் பொதுவாகக் கொள்வோம். எடுத்துக்காட்டாக, $s_1$,
  $s_2$,~$s_3$ என்ற தொடர்முறை~$C$ இல் இருந்தால், $s_2$, $s_3$ மற்றும்~$s_3$
  ஆகியவையும் அதில் இருக்கும். மேலும், இரு நிலைகள் $s_i$,~$s_j$ ஆகியவை ஒரு
  தொடர்முறை~$\sigma$ இல் தோன்றும்போது, $i < j$ எனில் $s_i \prec_\sigma s_j$
  என்கிறோம். $t \in V(p)$ எனில் $p$ ஆனது~$t$ இல் \emph{மெய்} என்கிறோம்.
\end{defn}

மாதிரி~$\mModel M$ இல் ஒரு காலப் புள்ளி~$t$ இல் !!a{formula}வின் மெய்மையை
மதிப்பிடும்போது, $t$ ஒரு நிலையாகத் தோன்றும் வரலாறு~$\sigma$ ஐச் சார்ந்தே
அதைச் செய்கிறோம் என்பதே தொடர்புடைய ஒரே மாற்றம். ஆனால் !!{propositional
variable}s அல்லது மெய்ச்சார்பு இணைப்பிகளுக்கான பொருண்மையியல் எதையும் நாம்
மாற்ற வேண்டியதில்லை. அவற்றில் எதுவும்~$\sigma$ ஐக் குறிப்பிடாது என்பதால்,
அவை அனைத்தும் \olref[sem]{defn:tmodels} இல் இருந்தவாறே உள்ளன. எனினும்,
நமது எதிர்காலச் செயற்குறி~$\Ftemp$ ஐ இப்போது மீள்வரையறுத்து,
இவ்வரலாறுகளைப் பொறுத்து நமது~$\Diamond$ செயற்குறியையும் சேர்க்கிறோம்.

\begin{defn}\ollabel{defn:phmodels}
  $\mModel M$ இல் \emph{$t, \sigma$ இல் !!a{formula}~$!A$ இன் மெய்மை},
  குறியீட்டில்: $\mSat{M}{!A}[t, \sigma]$:
  \begin{enumerate}
  \item\ollabel{defn:sub:phmodels-f} \indcase{!A}{\Ftemp
    !B}{$\mSat{M}{\indfrm}[t, \sigma]$ எனில், எனில் மட்டுமே,
    $t \prec_\sigma t'$ ஆகும் ஏதோவொரு $t' \in T$ க்கு
    $\mSat{M}{!B}[t', \sigma]$.}
  \item\ollabel{defn:sub:phmodels-diamond} \indcase{!A}{\Diamond
    !B}{$\mSat{M}{\indfrm}[t, \sigma]$ எனில், எனில் மட்டுமே, $t$ தோன்றும்
    ஏதோவொரு $\sigma' \in C$ க்கு $\mSat{M}{!B}[t, \sigma']$.}
  \end{enumerate}
\end{defn}

வேறு காலச் செயற்குறிகளையும் வாய்ப்புநிலைச் செயற்குறிகளையும் இதேபோல்
வரையறுக்கலாம். எனினும், வினைக்காலத்தையும் வாய்ப்புநிலையையும் இணைக்கும்
கூற்றுகளை இப்போது குறிக்க முடியும். எடுத்துக்காட்டாக, ``$p$~நிகழாது, ஆனால்
அது நிகழ்ந்திருக்கக் கூடியது'' என்பதை $\lnot \Ftemp p \land \Diamond \Ftemp p$
என்ற வாய்பாட்டால் குறிக்கலாம். ஒரு புள்ளி–வரலாறு இணையைத் தொடர்ந்து வரும்
எந்த நிலையிலும் $p$ மெய்யாகாமல், ஆனால் $p$ மெய்யாகும் ஓர் மாற்று வரலாறு
இருந்தால், அவ்வாய்பாடு அந்தப் புள்ளியிலும் வரலாற்றிலும் மெய்யாகும்.

\end{document}
